% Part: computability
% Chapter: recursive-functions
% Section: pr-relations

\documentclass[../../../include/open-logic-section]{subfiles}

\begin{document}

\olfileid[hi]{cmp}{rec}{prr}
\olsection{आदिम पुनरावर्ती संबंध}


\begin{defn}
किसी संबंध $R(\vec x)$ को आदिम पुनरावर्ती कहते हैं यदि उसका अभिलाक्षणिक
फलन
\[
\Char{R}(\vec x) = \left\{
  \begin{array}{ll}
  1 & \mbox{यदि $R(\vec x)$} \\
  0 & \mbox{अन्यथा}
  \end{array}
\right.
\]
आदिम पुनरावर्ती हो।
\end{defn}

दूसरे शब्दों में, जब हम किसी आदिम पुनरावर्ती संबंध $R(\vec x)$ की बात
करते हैं, तो हमारा आशय $\Char{R}(\vec
x) = 1$ रूप के संबंध से होता है,
जहाँ $\Char{R}$ ऐसा आदिम पुनरावर्ती फलन है जो किसी भी निवेश पर 1 या 0
लौटाता है। मसलन, संबंध $\fn{IsZero}(x)$, जो तभी और केवल तभी मान्य है जब
$x = 0$, उस फलन $\Char{\fn{IsZero}}$ के संगत है जिसे आदिम पुनरावर्तन से
इस प्रकार परिभाषित किया गया है:
\begin{align*}
\Char{\fn{IsZero}}(0) & = 1,\\
\Char{\fn{IsZero}}(x+1) & = 0.
\end{align*}

यह स्पष्ट होना चाहिए कि संबंधों का अन्य आदिम पुनरावर्ती फलनों के साथ
संयोजन किया जा सकता है। अतः निम्न भी आदिम पुनरावर्ती हैं:
\begin{enumerate}
\item समता संबंध $x = y$, जिसे $\fn{IsZero}(\left|x -
  y\right|)$ से परिभाषित किया गया है।
\item छोटा-या-बराबर संबंध $x \leq y$, जिसे $\fn{IsZero}(x
  \tsub y)$ से परिभाषित किया गया है।
\end{enumerate}


\begin{prop}
  आदिम पुनरावर्ती संबंधों का समुच्चय बूलीय संक्रियाओं के अधीन संवृत है;
  अर्थात् यदि $P(\vec x)$ और $Q(\vec x)$ आदिम पुनरावर्ती हैं, तो निम्न
  भी आदिम पुनरावर्ती हैं:
  \begin{enumerate}
  \item $\lnot P(\vec x)$
  \item $P(\vec x) \land Q(\vec x)$
  \item $P(\vec x) \lor Q(\vec x)$
  \item $P(\vec x) \lif Q(\vec x)$
  \end{enumerate}
\end{prop}

\begin{proof}
  मान लें कि $P(\vec x)$ और $Q(\vec x)$ आदिम पुनरावर्ती हैं, अर्थात्
  उनके अभिलाक्षणिक फलन $\Char{P}$ और $\Char{Q}$ ऐसे हैं। हमें दिखाना है
  कि $\lnot P(\vec x)$ इत्यादि के अभिलाक्षणिक फलन भी आदिम पुनरावर्ती हैं।
  \[
  \Char{\lnot P}(\vec x) = \begin{cases}
    0 & \text{यदि $\Char{P}(\vec x) = 1$}\\
    1 & \text{अन्यथा}
  \end{cases}
  \]
  हम $\Char{\lnot P}(\vec x)$ को $1 \tsub \Char{P}(\vec x)$ के रूप
  में परिभाषित कर सकते हैं।
  \[
  \Char{P \land Q}(\vec x) = \begin{cases}
    1 & \text{यदि $\Char{P}(\vec x) = \Char{Q}(\vec x) = 1$}\\
    0 & \text{अन्यथा}
  \end{cases}
  \]
  हम $\Char{P \land Q}(\vec x)$ को $\Char{P}(\vec x) \cdot
  \Char{Q}(\vec x)$ अथवा $\fn{min}(\Char{P}(\vec x), \Char{Q}(\vec
  x))$ के रूप में परिभाषित कर सकते हैं। इसी प्रकार,
  \begin{align*}
    \Char{P \lor Q}(\vec x) & = \fn{max}(\Char{P}(\vec x), \Char{Q}(\vec x)) \text{ और}\\
    \Char{P \lif Q}(\vec x) & = \fn{max}(1 \tsub
  \Char{P}(\vec x), \Char{Q}(\vec x)).
  \end{align*}
\end{proof}

\begin{prop}
  आदिम पुनरावर्ती संबंधों का समुच्चय परिबद्ध परिमाणीकरण के अधीन संवृत
  है; अर्थात् यदि $R(\vec x, z)$ कोई आदिम पुनरावर्ती संबंध है, तो निम्न
  संबंध भी ऐसे हैं:
  \begin{align*}
    & \bforall{z < y}{R(\vec x, z)} \text{ और}\\
    & \bexists{z < y}{R(\vec x, z)}.
  \end{align*}
  $\bforall{z < y}{R(\vec x, z)}$, $\vec x$ और $y$ के लिए तभी और केवल
  तभी मान्य है जब $R(\vec x, z)$ प्रत्येक~$z$ के लिए मान्य हो जो~$y$ से
  छोटा है; और $\bexists{z < y}{R(\vec x, z)}$ के लिए भी इसी प्रकार।
\end{prop}

\begin{proof}
  परिपाटी से हम $\bforall{z < 0}{R(\vec x, z)}$ को सत्य मानते हैं
  (इस सरल कारण से कि कोई $z$,~$0$ से छोटा नहीं है) और
  $\bexists{z < 0}{R(\vec x, z)}$ को असत्य मानते हैं। कोई परिबद्ध
  सार्विक परिमाणक ठीक परिमित गुणनफल या पुनरावृत्त न्यूनतम की तरह काम करता
  है; अर्थात् यदि $P(\vec x, y) \defiff \bforall{z <
  y}{R(\vec x, z)}$, तो $\Char{P}(\vec x, y)$ को इस प्रकार परिभाषित
  किया जा सकता है:
  \begin{align*}
    \Char{P}(\vec x, 0) & = 1\\
    \Char{P}(\vec x, y+1) & =
    \fn{min}(\Char{P}(\vec x, y), \Char{R}(\vec x, y)).
  \end{align*}
  परिबद्ध अस्तित्वात्मक परिमाणीकरण को इसी प्रकार $\fn{max}$ का उपयोग
  करके परिभाषित किया जा सकता है। वैकल्पिक रूप से, तुल्यता
  $\bexists{z < y}{R(\vec x, z)}
  \liff \lnot \bforall{z < y}{\lnot R(\vec x, z)}$ का उपयोग करके उसे
  परिबद्ध सार्विक परिमाणीकरण से परिभाषित किया जा सकता है। ध्यान दें कि,
  मसलन, $\bexists{x \leq y}{\dots
  x\dots}$ रूप का परिबद्ध परिमाणक
  $\bexists{x < y+1}{\dots x \dots}$ के तुल्य है।
\end{proof}

\begin{prob}
  दिखाएँ कि त्रि-स्थानी संबंध $x \equiv y \mod n$ (मापांक~$n$ में
  सर्वांगसमता) आदिम पुनरावर्ती है।
\end{prob}

एक अन्य उपयोगी आदिम पुनरावर्ती फलन सप्रतिबंध फलन $\fn{cond}(x,y,z)$ है,
जिसे इस प्रकार परिभाषित किया गया है:
\begin{align*}
  \fn{cond}(x,y,z) & = \begin{cases}
  y & \text{यदि $x = 0$} \\
  z & \text{अन्यथा}.
\end{cases}
\intertext{इसे पुनरावर्ती रूप से इस प्रकार परिभाषित किया गया है:}
\fn{cond}(0,y,z) & = y,\\
\fn{cond}(x+1,y,z) & = z.
\end{align*}
इसका उपयोग आदिम पुनरावर्ती संबंधों से आदिम पुनरावर्ती फलनों की
प्रकरणानुसार परिभाषाओं को उचित ठहराने के लिए किया जा सकता है:

\begin{prop}
यदि $g_0(\vec x)$, \dots,~$g_m(\vec x)$ आदिम पुनरावर्ती फलन हैं और
$R_0(\vec
x)$, \dots, $R_{m-1}(\vec x)$ आदिम पुनरावर्ती संबंध हैं, तो
निम्न प्रकार परिभाषित फलन $f$
\[
f(\vec x) = \begin{cases}
    g_0(\vec x) & \text{यदि $R_0(\vec{x})$} \\
    g_1(\vec x) & \text{यदि $R_1(\vec{x})$ और $R_0(\vec{x})$ नहीं} \\
    \vdots & \\
    g_{m-1}(\vec x) & \text{यदि $R_{m-1}(\vec{x})$ और पिछले में से कोई
      मान्य नहीं}
    \\
    g_m(\vec x) & \mbox{अन्यथा}
\end{cases}
\]
भी आदिम पुनरावर्ती है।
\end{prop}

\begin{proof}
  जब $m = 1$, तब यह केवल निम्न प्रकार परिभाषित फलन है:
  \[
  f(\vec x) = \fn{cond}(\Char{\lnot R_0}(\vec x),g_0(\vec x),g_1(\vec
  x)).
  \]
  $m$, $1$ से बड़ा होने पर इस रूप की परिभाषाओं का संयोजन भर करना होता है।
\end{proof}
\end{document}
